\documentclass[12pt]{article}
\usepackage[utf8]{inputenc}
\usepackage[T1]{fontenc}
\usepackage{hyperref}
\title{Urim Language Resources for K-12 Curriculum Development}
\author{Compiled by Claude Sonnet 5 (claude-sonnet-5) for Andr\'as Kornai}
\date{}
\begin{document}
\maketitle

\section{Introduction}
Urim (ISO 639-3 uri; also called Kalp) is a Torricelli language of East Sepik Province, Papua New Guinea. It has dialects Kukwo and Yangkolen. Speaker counts vary by source: about 3,700 (Wikipedia, 2003), about 4,000 (LingWeb) and 2,538 (Verbix). Do not confuse it with Urimo (urx), which is a different language.

\section{Primary Educational Resources}
None found in the search.

\section{Linguistic and Reference Materials}
\begin{itemize}
\item Wikipedia, Urim language. It cites the grammatical description by Hemmilä and Luoma (2009): \url{https://en.wikipedia.org/wiki/Urim_language}
\item Wikipedia, Torricelli languages (family context): \url{https://en.wikipedia.org/wiki/Torricelli_languages}
\item MPI-EVA channumerals page for Urim: \url{https://lingweb.eva.mpg.de/channumerals/Urim.htm}
\item WALS entry: \url{https://wals.info/languoid/lect/wals_code_uri}
\end{itemize}

\section{Religious and Community Resources}
None found.

\section{Digital and Audio Resources}
\begin{itemize}
\item Verbix endangered-languages page for Urim: \url{https://docs.verbix.com/EndangeredLanguages/Urim}
\end{itemize}

\section{Curriculum Development Context}
Urim is in daily use by the whole community but has no known school materials. The existing grammar (Hemmilä and Luoma 2009) is the basis for any orthography and teaching-material work. Only a grammar was found, with no dictionary or readers.
\end{document}